\chapter{Conclusions}
\label{conclusion}

% ── Key Findings ──────────────────────────────────────────────────────────────
\section{Key Findings}
\label{sec:summary}

This study empirically investigates the trade-off between predictive performance and ecological cost across classical models and deep \gls{nn}s on tabular data. By synthesizing the results of \expref{1} through \expref{5}, several core conclusions emerge that challenge conventional \gls{ml} practices.

\paragraph{Ecological efficiency is context-dependent}
Across the three datasets evaluated, no single model consistently led the efficiency rankings. The most ecologically viable choice shifted systematically with task scale and hardware alignment: \gls{lr} on Wine, \gls{xgb} \gls{cpu} on Credit, and \gls{gpu}-accelerated \gls{xgb} on HIGGS. This pattern confirms that ecological efficiency is not a fixed property intrinsic to any model, but emerges from the interaction between computational structure, hardware, and data volume. To properly quantify this shifting dynamic, this study introduced and validated the \gls{cf1} metric. \Gls{cf1} proved to be a highly effective absolute heuristic for exposing hidden ecological costs that traditional performance-only benchmark tables entirely obscure. On HIGGS, it penalized the \gls{mlp} despite its highest raw \gls{wf1} in the entire study, and on Credit it surfaced the twofold emissions gap between \gls{rf} and \gls{xgb} despite near-identical predictive performance. Through this \gls{cf1} lens, it becomes evident that efficiency emerges strictly from the fit between a model's computational profile, the hardware it runs on and the volume of data it processes. This is the empirical answer to \rqref{1}.

\paragraph{Hardware acceleration requires a minimum dataset scale}
The relationship between hardware and efficiency is sharply defined by dataset scale, challenging the assumption that \glspl{gpu} are universally beneficial. This study identified a hard empirical \gls{gpu} break-even point at approximately 80{,}000 rows for \gls{xgb} on tabular data. Below this threshold, \gls{gpu} acceleration acts as an ecological liability. On the Wine dataset, \gls{xgb} \gls{gpu} emitted 1.3 times the carbon emissions of its \gls{cpu} counterpart to achieve the exact same \gls{wf1}, because the fixed costs of device initialization and PCIe memory transfer completely dominated the actual computation. Once dataset scale surpasses the break-even point and saturates the parallel cores, the architecture pays off: at 11 million rows, the \gls{gpu} variant was seven times more emissions-efficient and nine times faster than the \gls{cpu} variant. The training-side advantage does not, however, survive into the deployment lifecycle. \gls{xgb} \gls{gpu}'s per-prediction inference cost of 40.71\,\textmu\gls{gco2eq} exceeds that of the \gls{cpu} variant (23.47\,\textmu\gls{gco2eq}) by a factor of 1.7, because every \gls{gpu} prediction call incurs a fixed device-dispatch overhead that large training batches fully amortize but individual inference calls cannot. The cumulative lifecycle emission curves of the two variants cross at approximately 137{,}000 predictions, the point at which the training footprint saved by \gls{gpu} acceleration has been fully consumed by the higher per-prediction cost. Beyond that volume, the \gls{cpu} variant is the more ecologically efficient choice for every additional query served. Hardware selection is therefore a two-dimensional decision: dataset scale determines whether \gls{gpu} acceleration is justified for training, and deployment volume determines whether that advantage is preserved over the full operational lifetime. These findings from \expref{1}, \expref{2}, and \expref{5} directly answer \rqref{3}.

\paragraph{Model structure outweighs hyperparameter tuning}
Beyond hardware, the structural design of a model dictates efficiency far more rigidly than hyperparameter tuning, particularly for tabular data. On the Credit dataset, \gls{rf} and \gls{xgb} \gls{cpu} reached a statistically indistinguishable \gls{wf1} (79.7\,\% versus 79.8\,\%), yet \gls{rf} incurred more than twice the emissions cost. No hyperparameter configuration can eliminate the inherent per-tree structural overhead of coarse-grained ensemble construction that causes this gap. At the scale of HIGGS, this overhead compounds to the point at which \gls{rf} had to be excluded from the largest dataset entirely, making it the only model in this benchmark that could not be evaluated on all three datasets and, by that criterion, the worst-scaling model in the study. When evaluating \gls{nn} architectures, the study found that width is the dominant axis for improving the \gls{mlp} on HIGGS, whereas depth primarily acts as computational overhead. Increasing layer width from 128 to 512 neurons improved the \gls{wf1} meaningfully, but adding depth beyond two layers yielded near-zero predictive return while monotonically increasing the emissions cost from 48\,\gls{gco2eq} to 93\,\gls{gco2eq}. These architectural findings from \expref{3} constitute the empirical answer to \rqref{2}.

\paragraph{Sublinear predictive returns impose a strict ecological boundary}
These structural limitations are further compounded by the sublinear scaling of predictive returns, proving that the dataset scaling curve acts as a strict ecological boundary. As data volume increases, models quickly hit capacity saturation or severe diminishing returns. \gls{lr}'s performance on HIGGS plateaued before 100{,}000 rows, rendering any compute expended on the remaining 8.7 million rows of the training pool ecologically wasted. Even for complex non-linear models capable of absorbing more data, expanding the HIGGS training data from one million to 8.8 million rows improved the \gls{mlp}'s \gls{wf1} by only 2.2 percentage points but exacted a 29-fold increase in emissions. This scaling pattern, documented in \expref{2}, is the empirical basis for answering \rqref{4}.

\paragraph{\gls{hpo} is a first-order ecological cost}
Focusing solely on the final training run fundamentally misrepresents the true ecological footprint of \gls{ml} systems. Hyperparameter tuning is a massive hidden first-order ecological cost. A deliberately constrained \texttt{Optuna} search for the \gls{mlp} on HIGGS emitted 3{,}180\,\gls{gco2eq} and took over 72 hours, making the initialization phase approximately 54 times costlier than the 58.5\,\gls{gco2eq} benchmark training run it ultimately enabled. Standard benchmark tables entirely omit this ecological initialization cost that precedes every deployment.

\paragraph{Deployment volume governs the lifecycle footprint}
The deployment lifecycle similarly overturns efficiency rankings based on inference latency. Because \gls{rf} is computationally heavy during inference, it breaks even at just 66 predictions on the Wine dataset, exhausting its modest training budget within minutes of active serving. In stark contrast, the \gls{mlp} on HIGGS features extremely fast inference and does not reach its break-even point until approximately 8.3 million predictions. The cross-model comparison between \gls{xgb} \gls{gpu} and \gls{mlp} on HIGGS makes the stakes concrete. Despite \gls{xgb} \gls{gpu} entering deployment with a training footprint 149 times smaller (0.393\,\gls{gco2eq} versus 58.5\,\gls{gco2eq}), its inference cost of 40.7\,µg per prediction is 5.8 times higher than \gls{mlp}'s 7.1\,µg. The two models reach equivalent total lifecycle emissions at approximately 1.73 million predictions. Beyond that point, \gls{mlp} is the more ecologically efficient choice for every additional query served. For high-volume deployments, inference latency per prediction entirely governs the lifecycle footprint, meaning that optimizing solely for training efficiency can lead to the wrong ecological choice. These findings from \expref{5} directly answer \rqref{5}.

\paragraph{Measurement accuracy and electricity grid timing are primary efficiency levers}
Finally, the accuracy of ecological benchmarking depends heavily on measurement methodology and external grid factors. Relying on software estimation tools like \texttt{CodeCarbon} can lead to systematic underestimation, particularly on \gls{cpu}-intensive workloads. Hardware-corrected measurements in this study exceeded raw \texttt{CodeCarbon} figures across every model-dataset combination, with a median underestimation factor of 1.92. This distortion is worst for short training runs, meaning published studies using uncorrected estimators likely report only a fraction of their true emissions. Even when accurately measured, absolute emissions remain highly volatile due to temporal grid variations. The carbon intensity of the German electricity grid fluctuates predictably, and the counterfactual analysis proves that grid-aware scheduling, for example running workloads during summer midday ($193\,\frac{\gls{gco2eq}}{\gls{kwh}}$) instead of winter evenings ($389\,\frac{\gls{gco2eq}}{\gls{kwh}}$), can reduce emissions by up to 70\,\% without altering a single line of code. For large-scale workloads, this zero-cost scheduling optimization saves more emissions than the combined training footprints of the smaller datasets, establishing timing as a primary lever for Green \gls{ai}, as quantified in \expref{4}.


% ── Implications for Green AI ─────────────────────────────────────────────────
\section{Implications for Green AI}
\label{sec:implications}

The empirical findings of this study, combined with the broader sustainability literature, provide several actionable implications for practitioners and researchers aiming to align \glspl{ml} development with Green \gls{ai} principles. Ecological efficiency requires moving away from unconditional defaults and adopting a holistic, lifecycle-aware approach to model design.

\paragraph{Data and model selection}
A fundamental reorientation is required in how datasets and model architectures are chosen. Dataset size must be treated as an aggressively tunable hyperparameter rather than a fixed asset to be unconditionally maximized. Because predictive returns on additional data generally follow a concave curve while energy costs grow linearly or faster \cite{koshuteRecommendingTrainingSet2021}, practitioners should identify the elbow point at which \gls{wf1} plateaus and halt data ingestion there. Assuming that all available data must be used is an ecological liability rather than a methodological virtue \cite{verdecchiaDataCentricGreenAI2022a}. Model selection must follow the same principle of matching computational profile to task scale. As the Credit results in \Cref{sec:results_credit} show, two structurally different models can reach the same \gls{wf1} while differing twofold in emissions, a gap that no amount of tuning can close \cite{schwartzGreenAI2020}. Unconditionally using a \gls{gpu} for small or medium tabular workloads compounds this problem further. The break-even analysis in this study demonstrates that below approximately 80{,}000 rows, device initialization overhead dominates the computation and increases rather than reduces emissions \cite{wuSustainableAIEnvironmental2022}. Even above that threshold, where \gls{gpu} training is ecologically justified, the advantage does not automatically extend to inference: per-prediction dispatch overhead that large training batches fully amortize cannot be amortized by individual prediction calls, as the lifecycle analysis in \Cref{sec:results_lifecycle} demonstrates for \gls{xgb}. Training on \gls{gpu} and serving predictions on \gls{cpu} is therefore a practical strategy that captures the training-side efficiency gains without incurring the inference-side overhead. For \gls{nn} architectures on tabular data specifically, excessive depth acts as a computational penalty rather than a predictive asset, since tabular datasets lack the strict local structure that makes deep hierarchies beneficial in image and text domains \cite{grinsztajnWhyTreebasedModels2022a}.

\paragraph{Lifecycle costs and deployment decisions}
Evaluating models exclusively on training efficiency obscures the true environmental cost of deployment. As inference can account for 80 to 90 percent of the total \gls{ml} energy demand in production \cite{desislavovTrendsAIInference2023}, ecological model selection must incorporate an explicit deployment volume estimate from the outset. The break-even prediction counts derived in \Cref{sec:results_lifecycle} make this concrete: an inference-heavy model can exhaust its training carbon budget within tens of predictions, whereas an inference-light one may not break even for millions. A training-efficient model that is inference-heavy is not a green choice at scale. Post-training optimization techniques such as model distillation, pruning, and quantization can substantially lower per-prediction energy without requiring retraining and should be considered a standard part of the deployment pipeline for \gls{nn}-based models such as the \gls{mlp} \cite{pattersonCarbonEmissionsLarge2021}. For tree-based models these techniques are largely inapplicable, and the primary inference-side efficiency lever there is model selection at training time, as the gap between \gls{rf} and \gls{xgb} inference costs in this study illustrates. Published efficiency comparisons must also account for the ecological cost of hyperparameter tuning, which for the \gls{mlp} on HIGGS dwarfed the final training run by the factor established in \Cref{sec:results_tuning}. Any benchmark that omits this overhead misrepresents the initialization cost of the models it compares \cite{strubellEnergyPolicyConsiderations2019a}. A further consideration is the rebound effect: as models become cheaper and more efficient to run, their usage often scales up, driving aggregate energy consumption higher. This dynamic, known as Jevons' Paradox, has been identified as a structural risk in the \gls{ai} efficiency literature \cite{luccioniPowerHungryProcessing2024}. This study does not directly observe or measure rebound effects. The mechanism is included as a theoretical implication drawn from that literature rather than an empirical result. The lifecycle framing developed here does, however, make the dynamic concrete: if the per-prediction efficiency improvements documented in \Cref{sec:results_lifecycle} were to lower the cost barrier and expand deployment volume proportionally, aggregate emissions could rise even as individual prediction costs fall. Efficiency gains must therefore be paired with deliberate usage constraints to produce net ecological benefits.

\paragraph{Measurement and infrastructure}
Accurate ecological accounting requires rigor at both the measurement and infrastructure levels. Software-based estimation tools that rely on \gls{tdp} values consistently undercount actual energy usage. The median hardware-corrected factor of 1.92 observed in this study implies that benchmark figures produced with uncorrected \texttt{CodeCarbon} on comparable hardware capture roughly 52\,\% of actual emissions. Direct hardware sensor readings should be the standard for any work making quantitative ecological claims \cite{rodriguezEvaluatingEnergyConsumption2024}. Once measurement is reliable, infrastructure timing becomes a powerful zero-cost lever. Shifting training jobs temporally to align with low-intensity grid hours, when solar or wind generation peaks, can reduce absolute emissions substantially without altering a single line of code. Beyond coarse rescheduling, carbon-aware training systems can monitor real-time carbon intensity over the course of a run and dynamically throttle or pause \gls{gpu} compute whenever the electricity grid mix turns carbon-heavy, lowering the training footprint while keeping the runtime penalty small \cite{yangChasingLowCarbonElectricity2023}. For large-scale workloads, the counterfactual analysis in this study shows this timing difference can exceed the combined training footprints of smaller datasets, making carbon-aware scheduling one of the most cost-effective actions available to practitioners on renewable-heavy grids \cite{dodgeMeasuringCarbonIntensity2022}. In cloud environments, spatial shifting is often even more powerful. Because model training is rarely latency-bound, workloads can be routed to data centers in regions with predominantly renewable energy supply. The difference between a fossil-fuel-dependent and a hydro-heavy region can reduce emissions by a factor of up to 30 \cite{hendersonSystematicReportingEnergy2020}. Operational optimization must nevertheless be paired with an acknowledgment of embodied carbon. The manufacturing of specialized \gls{ai} hardware, including raw material extraction and the substantial water footprint of semiconductor fabrication, is increasingly recognized as a dominant factor in the overall environmental impact of \gls{ai} systems \cite{eilamMethodologyFrameworkAI2023,wuSustainableAIEnvironmental2022}. Green \gls{ai} therefore requires not only efficient use of existing hardware but also deliberate decisions about hardware procurement and lifecycle extension.

\paragraph{Systemic and community-level implications}
Long-term sustainability in \gls{ai} ultimately requires systemic shifts in how the research community structures its incentives and tools. Because many practitioners rely on default parameters in \gls{ml} libraries, shifting those defaults towards more energy-efficient configurations would reduce aggregate emissions across the entire community with minimal effort per user \cite{hendersonSystematicReportingEnergy2020}. The community should additionally adopt standardized reporting frameworks and energy leaderboards that track computational and carbon costs alongside traditional performance metrics \cite{pattersonCarbonEmissionsLarge2021}. Benchmark tables that list only performance figures reward parameter count and training time while rendering ecological cost invisible. By aligning academic and industrial incentives to recognize algorithmic efficiency as a first-class evaluation criterion alongside raw predictive performance, the field can foster a virtuous cycle of innovation that does not treat environmental impact as an afterthought.



% ── Future Work ───────────────────────────────────────────────────────────────
\section{Future Work}
\label{sec:future_work}

Many of the most productive directions opened by this work follow directly from the boundaries within which it was conducted. The study's constraints are best read not as dead ends but as a map of where the next experiments are most needed, and each is paired below with the research direction that would resolve it. The directions fall into four concentric rings: methodological refinement, scope expansion, model extensions, and a systems-level synthesis that would transform ecological benchmarking from a post-hoc activity into a design discipline.

\subsection{Measurement and Reproducibility}

\paragraph{Hardware and electricity grid context}
Every absolute emission figure and break-even threshold in this study was produced on one machine connected to the German electricity grid, so the relative efficiency rankings are best treated as hypotheses for other hardware contexts rather than universal rules. The \gls{gpu} break-even threshold of approximately 80{,}000 rows, for instance, is a direct function of this \gls{gpu}'s initialization overhead and PCIe transfer cost, and a datacenter-class card with higher parallelism and faster interconnect would place it elsewhere, possibly by a large margin. This dependence is especially relevant for the cloud deployments that dominate production \gls{ml} infrastructure, where \gls{gpu} instances run on shared hardware billed by the second, their initialization overhead and per-compute carbon intensity differ structurally from a locally attached consumer card, and the host data center's carbon intensity can diverge from the German grid by an order of magnitude \cite{hendersonSystematicReportingEnergy2020}. Reproducing the benchmark on laptop-class \glspl{cpu}, datacenter-grade \glspl{gpu}, and low-power inference chips, and across grids with fundamentally different carbon profiles such as Nordic hydro-heavy or U.S.\ Midwest coal-heavy mixes \cite{dodgeMeasuringCarbonIntensity2022}, would reveal which findings are hardware-specific and which reflect structural properties of the models that generalize. Establishing the family of curves to which the break-even threshold belongs would give practitioners a directly applicable decision rule rather than a single data point, and it is a precondition for these conclusions to inform cloud procurement decisions. Wright et al.\ \parencite{wrightEfficiencyNotEnough2025} formalize this hardware-context sensitivity as a fundamental discrepancy between compute efficiency, energy efficiency, and carbon efficiency, three quantities that diverge most sharply when hardware utilization is misaligned with workload scale. That is precisely the condition this study documents for \gls{gpu} acceleration below the break-even threshold, and verifying whether the discrepancy persists on datacenter-class \glspl{gpu}, where initialization overhead and interconnect costs operate at a different scale, would determine how far the single consumer-hardware data point obtained here generalizes.

\paragraph{Platform independence}
Reproducibility also requires solving the platform dependency that currently limits the correction pipeline. The hardware-corrected measurements presented here rely on LibreHardwareMonitor, a Windows-only tool requiring administrator privileges. Linux exposes equivalent processor energy counters through the \gls{rapl} interface without such requirements, and reimplementing the same correction against \gls{rapl} would remove the platform constraint, allow independent cross-validation of the 1.92 underestimation factor on a second hardware context, and enable the methodology to travel to cloud-based research infrastructure where Windows instances are rarely available.

\paragraph{Measurement scope and asymmetries}
The accounting boundary of this study leaves several quantities unmeasured, and closing each of them would tighten the figures the benchmark reports. Measured against the Greenhouse Gas Protocol, which distinguishes Scope 1, Scope 2, and Scope 3 emissions \cite{dodgeMeasuringCarbonIntensity2022,liMakingAILess2025}, the present figures capture only part of Scope 2. \gls{ram} energy is estimated from the number of installed DIMM slots at a fixed wattage rather than measured directly, and data loading, preprocessing, and the train-test split are deliberately excluded from the tracker so that reported figures reflect training cost alone. Extending the tracker to instrument these phases, and reading memory power from hardware counters rather than a per-slot heuristic, would yield the true end-to-end footprint of a run, which on HIGGS is non-trivially higher than the training-only figure suggests. Three measurement asymmetries further limit cross-model comparability and each points to a concrete refinement. Because \gls{lr}'s single-threaded \gls{sag} solver runs at roughly 10\,\% \gls{cpu} load while \gls{rf} and \gls{xgb} saturate all eight cores, the static idle component of Package Power inflates \gls{lr}'s measured emissions, so subtracting a measured idle baseline would isolate the algorithmic cost more cleanly. Because both \gls{xgb} \gls{gpu} and \gls{mlp} route prediction through the \gls{gpu}, the \texttt{CPUPowerMonitor} captures only their host-side overhead, leaving their reported per-prediction costs and break-even counts as lower bounds. Instrumenting \gls{gpu}-side inference energy directly through an interface such as \gls{nvml} would replace those lower bounds with true values and sharpen the lifecycle break-even estimates. Cleanly separating idle draw from active workload power remains a recognized open problem in energy measurement for \gls{ml} systems more broadly \cite{tschandMLPerfPowerBenchmarking2025}.

\paragraph{Statistical robustness}
Every model-dataset combination was evaluated once on a single fixed 80/20 split under one random seed, so no confidence intervals or significance tests accompany the reported figures, and differences as fine as \gls{xgb} \gls{cpu}'s \gls{wf1} of 79.8\,\% against \gls{rf}'s 79.7\,\% on Credit rest on structural reasoning rather than hypothesis testing. Repeated runs across multiple seeds or a repeated k-fold evaluation would convert each point estimate into a distribution and allow genuine variance to be distinguished from systematic difference. This refinement matters most for Wine, where the roughly 36 test samples yield an estimate too coarse to separate models, all five of which returned an identical \gls{wf1} of 97.2\,\%. The scaling experiment already replicates every model-subset combination under three independent seeds and reports the mean, so reporting the variance around those means would be a small extension that communicates the remaining uncertainty at small subset sizes, where a single atypical subsample draw can shift the result.

% LTeX: enabled=false
\paragraph{Dynamic carbon intensity}
A further refinement concerns the carbon intensity inputs themselves. The counterfactual scheduling analysis in \Cref{sec:results_carbon} relies on average hourly electricity grid carbon intensity for Germany. Consequential carbon accounting argues that the relevant quantity is dynamic carbon intensity: the emissions rate of the generation unit that would actually be dispatched in response to an additional unit of demand \cite{dodgeMeasuringCarbonIntensity2022}. Dynamic carbon intensity is more volatile than the average and frequently higher during demand peaks when fossil backup capacity is activated. Rerunning the scheduling analysis with dynamic carbon intensity data would test whether the 70 percent emissions reduction identified under average carbon intensity assumptions survives a more demanding accounting standard, and it would yield a more precise estimate of the true climate impact of carbon-aware scheduling decisions. The same time resolution would address a related boundary in the present figures: because \texttt{CodeCarbon} converts measured energy with a fixed annual average of $381\,\frac{\gls{gco2eq}}{\gls{kwh}}$ for Germany, absolute \gls{cf1} values cannot be compared directly across datasets whose runs span very different durations and therefore different exposure to grid fluctuation. Converting energy with the carbon intensity recorded at each run's actual execution time would remove that distortion.
% LTeX: enabled=true

\subsection{Model and Domain Scope}

\paragraph{Gradient-boosting variants}
Once the measurement foundation is more robust, extending the model landscape is most valuable when each addition is tied to a specific structural hypothesis rather than used purely to increase coverage. The most targeted extension is a direct comparison of \gls{xgb} against LightGBM \cite{keLightGBMHighlyEfficient2017} and CatBoost \cite{prokhorenkovaCatBoostUnbiasedBoosting2018}: both operate within the gradient-boosting family but use leaf-wise growth and native categorical encoding, structural choices that alter memory access patterns and core utilization in ways that could shift the efficiency profile relative to \gls{xgb}'s level-wise algorithm. This would test whether the efficiency gap between gradient-boosting and ensemble methods is a property of the family or of \gls{xgb}'s specific implementation.

\paragraph{Attention-based tabular architectures}
A second targeted opportunity is offered by attention based tabular architectures such as TabNet \cite{arikTabNetAttentiveInterpretable2021} and FT-Transformer \cite{gorishniyRevisitingDeepLearning2021a}: their sequential feature selection and self-attention mechanisms create fundamentally different compute graphs from tree-based and dense-layer models, and their per-prediction latency and scaling behavior under the lifecycle break-even framework could yield results that neither family analyzed here would anticipate.

\paragraph{Beyond tabular data}
All experiments here cover tabular classification on three datasets chosen as a deliberate size spectrum rather than a broad sample of tabular problems, so the conclusions about depth, width, hardware selection, and dataset scaling do not transfer without re-evaluation to computer vision, \gls{nlp}, or regression, where architectures map onto hardware in fundamentally different ways and different feature dimensionalities, class structures, or noise profiles could alter the rankings observed here. The most direct scope expansion is therefore an evaluation of the same benchmarking framework on image classification and natural language tasks. The depth-as-overhead finding in this study is structurally conditioned on the absence of local spatial or sequential feature dependencies in tabular data. In convolutional and transformer-based architectures applied to images and text, hierarchical depth is not overhead but the primary mechanism by which useful representations are formed: removing layers degrades performance in a way that has no analog in the tabular \gls{mlp} experiments reported here \cite{grinsztajnWhyTreebasedModels2022a}. Applying the \gls{cf1} framework to a ResNet or \gls{bert} variant would therefore test whether depth remains a net ecological liability once it carries genuine predictive weight, or whether the relationship inverts entirely. The \gls{gpu} break-even analysis is equally likely to shift in these domains. Vision and language models are typically an order of magnitude larger than the tabular classifiers studied here \cite{pattersonCarbonEmissionsLarge2021}, and their training procedures are designed around \gls{gpu}-native tensor operations. The fixed initialization overhead that constitutes an ecological penalty at small tabular scale may be negligible against the compute budget of a full pretraining run. % LTeX: enabled=false
Replicating the lifecycle and scaling analyzes in these settings would reveal which efficiency conclusions reflect structural properties of the models and which are artifacts of the tabular domain.
% LTeX: enabled=true

\subsection{Model Extensions}

\paragraph{Carbon-aware \gls{hpo}}
This study found that a deliberately constrained \texttt{Optuna} search for the \gls{mlp} on HIGGS emitted approximately 54 times the final training run's footprint, and that number reflects a search space kept narrow by design. Two deliberate trade-offs shape the reported figures and define a clean agenda for this direction. First, the trial budget was held uniform across all models so that search effort stays symmetric, which means the \gls{mlp}'s conditional space of effective dimensionality five to eight is harder to map in limited trials than the low-dimensional spaces of \gls{rf} or \gls{xgb}, leaving reported \gls{mlp} performance a lower bound on what a more exhaustive search could reach. Second, \gls{lr} was left untuned to represent the zero-tuning-cost end of the spectrum, so part of its \gls{wf1} gap to the tuned gradient-boosting models reflects the contrast between an untuned linear model and optimized non-linear ones, and a regularization search could close some of it, particularly on Credit. The \gls{mlp} also omits \gls{amp} training for hardware-independent comparability, so its reported emissions are an upper bound on a deployment-grade implementation. Incorporating \gls{co2eq} or total energy as a secondary objective within the \gls{hpo} loop would allow practitioners to navigate the \gls{wf1}-emissions trade-off during model selection rather than discovering it after the optimization budget has already been spent. Paired with early stopping on per-trial energy expenditure and a tuning-inclusive version of \gls{cf1} as the optimization target, carbon-aware \gls{hpo} could substantially reduce the ecological initialization cost of large-scale training without sacrificing competitive \gls{wf1}. Wright et al.\ \parencite{wrightEfficiencyNotEnough2025} identify, however, a behavioral dimension to this problem: practitioners given a more efficient training pipeline frequently respond by running larger hyperparameter searches rather than banking the savings, effectively reinstating the original carbon budget. Measuring whether that behavioral rebound occurs in practice, for instance by tracking whether a twofold reduction in per-trial cost leads to a proportionally larger trial count, would establish whether carbon-aware \gls{hpo} produces a genuine net reduction or whether the efficiency gain is absorbed by expanded search.

\paragraph{Data-efficient learning}
The complementary problem lies on the data side. This study identifies the scaling elbow at which additional data yields diminishing predictive returns relative to \gls{co2eq} cost but does not examine how to extract more signal from fewer examples once that elbow is known. Active learning, core-set selection, and data distillation methods could potentially shift the \gls{wf1} curve above the observed saturation point without expanding the training set, and evaluating them on HIGGS subsets at the saturation points documented in \Cref{sec:results_scaling} would quantify whether the elbow is a movable threshold or a model capacity ceiling. The scaling experiment that established this elbow carried hyperparameters over from the full 11-million-row tuning run rather than re-optimizing per subset, which isolates the pure effect of data volume but renders the reported \gls{wf1} figures at 1{,}000 or 10{,}000 instances conservative lower bounds for models configured for a much larger dataset. \gls{rf} was likewise excluded from HIGGS runs above 500{,}000 rows because the per-run cost at full scale was ecologically unjustifiable for the insight it would have added. Re-optimizing per subset under an explicit carbon budget, and extending the curve for the heavier ensemble models, would refine the elbow into a per-model decision rule.

\subsection{Systems-Level Synthesis}

\paragraph{Total carbon ownership framework}
The most ambitious direction is to move from isolated efficiency measurement toward what Wright et al.\ \parencite{wrightEfficiencyNotEnough2025} describe as systems thinking in Green \gls{ai}: constructing a total carbon ownership framework that models the causal interactions between all identified efficiency levers simultaneously rather than analyzing them in sequence after the fact. Model selection, hardware configuration, dataset size, tuning budget, temporal electricity grid scheduling, and deployment volume were each analyzed in isolation throughout this work. A unified framework would map these dimensions jointly: given a dataset, a deployment traffic estimate, and a target \gls{wf1}, it would project a full lifecycle \gls{co2eq} budget across every pipeline phase before the first training run begins. That transition from post-hoc reporting to pre-deployment carbon budgeting would represent a genuine shift in how ecological impact is addressed in \glspl{ml}.

\paragraph{Inference optimization}
The inference side of the lifecycle remains equally open. Post-training optimization techniques such as quantization, pruning, and knowledge distillation reduce per-prediction compute without requiring a full retraining cycle \cite{pattersonCarbonEmissionsLarge2021}, and measuring the \gls{co2eq} cost of those compression steps alongside the inference speedup they provide would extend the lifecycle model developed in \Cref{sec:results_lifecycle} with a compression phase, potentially shifting the break-even prediction counts significantly for the heavier models.

\paragraph{Embodied carbon}
A complete lifecycle framework must also address the embodied carbon dimension that Wright et al.\ \parencite{wrightEfficiencyNotEnough2025} identify as a third structural discrepancy: the gap between operational emissions and the manufacturing footprint of the underlying hardware. This study accounts for no embodied carbon at all, excluding the manufacturing emissions of the \gls{cpu}, \gls{gpu}, and \gls{ram} entirely, a gap that bites hardest for the short training runs on small datasets where that component can outweigh the operational footprint and therefore narrows what the operational figures alone can establish. Amortizing the hardware manufacturing footprint across training runs and inference volume would place procurement decisions on the same carbon ledger as model and hardware choices, and could fundamentally alter break-even thresholds for hardware-intensive models on short-horizon deployments.

% ── Concluding Remarks ────────────────────────────────────────────────────────
\section{Concluding Remarks}
\label{sec:concluding_remarks}

This work set out to ask a deceptively simple question: which model is the most ecologically efficient choice for tabular classification, and under what conditions does that answer change? The answer that emerged is clear enough to be actionable but complex enough to resist any single formula. Ecological efficiency is not an intrinsic property of a model. It emerges from the fit between a model's computational structure, the hardware it runs on, the volume of data it processes, the electricity grid it draws power from, and the scale at which it will eventually serve predictions. No benchmark table that reports performance alone can capture this. The contribution of this study is not a ranked list of preferred classifiers. It is a set of empirical instruments, a corrected measurement methodology, a context-sensitive efficiency metric, and a lifecycle accounting framework, that together make the question answerable in a given deployment context rather than in the abstract. Green \gls{ai} does not require a fundamentally different kind of engineering. It requires the same rigorous, constraint-aware reasoning that good engineering has always demanded, applied consistently to a dimension that the field has for too long treated as external to the problem.